% GENERATED FILE. Edit canonical lessons, then run npm run generate:books.
% canonical-source-hash: fnv1a64:dfb7208397a5127d
% canonical-lessons: RU-C243-molodoy, RU-C243-pozhiloy, RU-C243-vysokiy, RU-C243-nizkiy, RU-C243-nekrasivyy

\chapter{Young, Elderly, Tall, Low, Ugly}
\label{ch:ru-young-elderly-tall-low-ugly}
\hlchaptermodality{243}

\begin{chapteropening}
\textbf{By the end of this chapter:} \emph{I can say young, elderly, tall, low and ugly.}

Complete the last lesson of chapter 243: I can say young, elderly, tall, low and ugly.
\end{chapteropening}

\section[\texorpdfstring{\ru{молодой} (molodóy)}{молодой (molodóy)}]{\ru{молодой} (molodóy) — young}

\label{lesson:RU-C243-molodoy}



\begin{quote}
Before the new one: say the Russian for a zoo, then the Russian for a circus.
\end{quote}

\subsection*{You'll want to know: \ru{молодой}}
\textbf{\ru{молодой}} (\emph{molodóy}) — \textquotedblleft{}young\textquotedblright{}.

\begin{grammarlens}[title={What you've built}]
One more word for this chapter.
\end{grammarlens}

\subsection*{Your turn}
\begin{itemize}
\raggedright
  \item \emph{Say it:} \emph{molodóy}
  \item \emph{Say it:} \emph{molodóy}, once more
  \item \emph{Say it:} say \emph{molodóy}
  \item \emph{Recall:} say the Russian for a zoo, then the Russian for a circus, then say \emph{molodóy} again
\end{itemize}

\begin{tcolorbox}[breakable,colback=teal!4,colframe=teal!35!black,title={Before you move on}]
What does \ru{молодой} mean? (\textquotedblleft{}Young\textquotedblright{}.) Say it once more.
\end{tcolorbox}
\section[\texorpdfstring{\ru{пожилой} (pozhilóy)}{пожилой (pozhilóy)}]{\ru{пожилой} (pozhilóy) — elderly}

\label{lesson:RU-C243-pozhiloy}



\begin{quote}
Before the new one: say the Russian for a circus, then the Russian for young.
\end{quote}

\subsection*{You'll want to know: \ru{пожилой}}
\textbf{\ru{пожилой}} (\emph{pozhilóy}) — \textquotedblleft{}elderly\textquotedblright{}.

\begin{grammarlens}[title={What you've built}]
One more word for this chapter.
\end{grammarlens}

\subsection*{Your turn}
\begin{itemize}
\raggedright
  \item \emph{Say it:} \emph{pozhilóy}
  \item \emph{Say it:} \emph{pozhilóy}, once more
  \item \emph{Say it:} say \emph{pozhilóy}
  \item \emph{Recall:} say the Russian for a circus, then the Russian for young, then say \emph{pozhilóy} again
\end{itemize}

\begin{tcolorbox}[breakable,colback=teal!4,colframe=teal!35!black,title={Before you move on}]
What does \ru{пожилой} mean? (\textquotedblleft{}Elderly\textquotedblright{}.) Say it once more.
\end{tcolorbox}
\section[\texorpdfstring{\ru{высокий} (vysókiy)}{высокий (vysókiy)}]{\ru{высокий} (vysókiy) — tall, high}

\label{lesson:RU-C243-vysokiy}



\begin{quote}
Before the new one: say the Russian for young, then the Russian for elderly.
\end{quote}

\subsection*{You'll want to know: \ru{высокий}}
\textbf{\ru{высокий}} (\emph{vysókiy}) — \textquotedblleft{}tall, high\textquotedblright{}.

\begin{grammarlens}[title={What you've built}]
One more word for this chapter.
\end{grammarlens}

\subsection*{Your turn}
\begin{itemize}
\raggedright
  \item \emph{Say it:} \emph{vysókiy}
  \item \emph{Say it:} \emph{vysókiy}, once more
  \item \emph{Say it:} say \emph{vysókiy}
  \item \emph{Recall:} say the Russian for young, then the Russian for elderly, then say \emph{vysókiy} again
\end{itemize}

\begin{tcolorbox}[breakable,colback=teal!4,colframe=teal!35!black,title={Before you move on}]
What does \ru{высокий} mean? (\textquotedblleft{}Tall, high\textquotedblright{}.) Say it once more.
\end{tcolorbox}
\section[\texorpdfstring{\ru{низкий} (nízkiy)}{низкий (nízkiy)}]{\ru{низкий} (nízkiy) — low, short}

\label{lesson:RU-C243-nizkiy}



\begin{quote}
Before the new one: say the Russian for elderly, then the Russian for tall.
\end{quote}

\subsection*{You'll want to know: \ru{низкий}}
\textbf{\ru{низкий}} (\emph{nízkiy}) — \textquotedblleft{}low, short\textquotedblright{}.

\begin{grammarlens}[title={What you've built}]
One more word for this chapter.
\end{grammarlens}

\subsection*{Your turn}
\begin{itemize}
\raggedright
  \item \emph{Say it:} \emph{nízkiy}
  \item \emph{Say it:} \emph{nízkiy}, once more
  \item \emph{Say it:} say \emph{nízkiy}
  \item \emph{Recall:} say the Russian for elderly, then the Russian for tall, then say \emph{nízkiy} again
\end{itemize}

\begin{tcolorbox}[breakable,colback=teal!4,colframe=teal!35!black,title={Before you move on}]
What does \ru{низкий} mean? (\textquotedblleft{}Low, short\textquotedblright{}.) Say it once more.
\end{tcolorbox}
\section[\texorpdfstring{\ru{некрасивый} (nekrasívyy)}{некрасивый (nekrasívyy)}]{\ru{некрасивый} (nekrasívyy) — ugly}

\label{lesson:RU-C243-nekrasivyy}



\begin{quote}
Before the new one: say the Russian for tall, then the Russian for low.
\end{quote}

\subsection*{You'll want to know: \ru{некрасивый}}
\textbf{\ru{некрасивый}} (\emph{nekrasívyy}) — \textquotedblleft{}ugly\textquotedblright{}.

\begin{grammarlens}[title={What you've built}]
One more word for this chapter.
\end{grammarlens}

\subsection*{Your turn}
\begin{itemize}
\raggedright
  \item \emph{Say it:} \emph{nekrasívyy}
  \item \emph{Say it:} \emph{nekrasívyy}, once more
  \item \emph{Say it:} say \emph{nekrasívyy}
  \item \emph{Recall:} say the Russian for tall, then the Russian for low, then say \emph{nekrasívyy} again
\end{itemize}

\begin{tcolorbox}[breakable,colback=teal!4,colframe=teal!35!black,title={Before you move on}]
What does \ru{некрасивый} mean? (\textquotedblleft{}Ugly\textquotedblright{}.) Say it once more.
\end{tcolorbox}
